\begin{frame}{사전확률 $P(y)$ 의 힘: 그레이엄의 ``스팸이 90\%''}
\begin{center}\begin{varwidth}{0.92\textwidth}\usebeamercolor[fg]{normal text}
  \textbf{NB 점수 첫 항 $\log P(y)$ 가 사전확률}
  \vskip0.4em
  
  \[ P(\text{공짜}|\text{spam})\,P(\text{회의}|\text{spam})
     = \tfrac34\times\tfrac12 = \tfrac38 = \tfrac14\times\tfrac34
     = P(\text{공짜}|\text{ham})\,P(\text{회의}|\text{ham}) \]
  \vskip0.2em
  \textbf{가능도비 2 $\Rightarrow$ 분류는 $P(y)$ 로 결정}
\end{varwidth}\end{center}
\note{\begin{itemize}
\item NB 점수 $\log P(y) + \sum_j \log P(x_j|y)$ 의 \textbf{첫 항}
  $\log P(y)$ 가 바로 사전확률 --- 기저율 오류가 그대로 적용.
\item 앞의 학습 데이터로 만든 NB에서, \textbf{가능도가 전혀 도움이 안
  되는} 메일 $\{$공짜, 회의$\}$ 를 분류:
\item 가능도비는 정확히 \textbf{2}(스팸 쪽) --- ``스팸도 정상도 똑같이
  잘 설명'' $\Rightarrow$ 분류는 \textbf{완전히 $P(y)$ 로 결정}.
\end{itemize}}
\end{frame}
